% Part: first-order-logic
% Chapter: model-theory
% Section: partial-iso

\documentclass[../../../include/open-logic-section]{subfiles}

\begin{document}

\olfileid{mod}{bas}{dlo}
\section{सघन रेषीय क्रम}

\begin{defn}
  \emph{अंतबिंदुविरहित सघन रेषीय क्रम} म्हणजे एकच 2-स्थानी
  !!{predicate}~$<$ असलेल्या !!{language} साठीची आणि पुढील वाक्ये पूर्ण
  करणारी !!a{structure}~$\Struct{M}$ होय:
  \begin{enumerate}
  \item $\lforall[x][\lnot x < x]$;
  \item $\lforall[x][\lforall[y][\lforall[z][(x < y \lif (y < z \lif x
    <z ))]]]$;
  \item $\lforall[x][\lforall[y][(x< y \lor \eq[x][y] \lor y < x)]]$;
  \item $\lforall[x][\lexists[y][x < y]]$;
  \item $\lforall[x][\lexists[y][y < x]]$;
  \item $\lforall[x][\lforall[y][(x < y \lif \lexists[z][(x < z \land
        z < y))]]]$.
 \end{enumerate}
\end{defn}

\begin{thm}\ollabel{thm:cantorQ}
  कोणतेही दोन !!{enumerable} अंतबिंदुविरहित सघन रेषीय क्रम समरूपी
  असतात.
\end{thm}

\begin{proof}
  $\Struct{M_1}$ आणि $\Struct{M_2}$ हे !!{enumerable} अंतबिंदुविरहित
  सघन रेषीय क्रम असू द्या, जेथे ${<_1} = \Assign{<}{M_1}$ आणि ${<_2} =
  \Assign{<}{M_2}$; तसेच $\PIso{I}$ हा त्यांच्यामधील सर्व आंशिक
  समरूपणांचा संच असू द्या. किमान $\emptyset \in \PIso{I}$ असल्यामुळे
  $\PIso{I}$ रिक्त नाही. $\PIso{I}$ हा पुढे-मागे गुणधर्म पूर्ण करतो हे
  आपण दाखवतो. मग $\Struct{M_1} \iso[p] \Struct{M_2}$; आणि
  \olref[pis]{thm:p-isom1} नुसार प्रमेय सिद्ध होते.

  $\PIso{I}$ हा पुढे गुणधर्म पूर्ण करतो हे दाखवण्यासाठी $p \in
  \PIso{I}$ असू द्या, $i = 1$, \dots,~$n$ साठी $p(a_i) = b_i$ असू द्या
  आणि सर्वसाधारणतेला बाधा न आणता $a_1 <_1 a_2 <_1 \cdots <_1 a_n$
  असे समजा. $a \in \Domain{M_1}$ दिलेला असताना, तो आधीच त्या नकाशाच्या
  प्रांतात असेल, तर तोच नकाशा अपेक्षित विस्तार आहे. नकाशा रिक्त असेल,
  तर दुसऱ्या रचनेच्या अरिक्त प्रांतातील कोणताही घटक प्रतिमा म्हणून
  निवडून त्याची दिलेल्या घटकाबरोबरची जोडी त्या नकाशात घालता येते.
  उरलेल्या प्रकरणांत
  पुढीलप्रमाणे $b \in \Domain{M_2}$ शोधा:
  %READERNOTE{OLFOL-028 — गोठवलेल्या इंग्रजी पुराव्यात p च्या प्रांतात a आधीच असण्याचे आणि p रिक्त असण्याचे प्रसंग वगळले आहेत; परंतु I मध्ये रिक्त नकाशा स्पष्टपणे समाविष्ट आहे. मराठी पुराव्यात सूत्रे न बदलता हे दोन प्रसंग जोडले आहेत.}
  \begin{enumerate}
  \item जर $a <_1 a_1$, तर $b <_2 b_1$ असेल असा $b \in
    \Domain{M_2}$ घ्या;
  \item जर $a_n <_1 a$, तर $b_n <_2 b$ असेल असा $b \in \Domain{M_2}$ घ्या;
 \item एखाद्या $i$ साठी $a_i <_1 a <_1 a_{i+1}$ असेल, तर $b_i <_2 b
   <_2 b_{i+1}$ असेल असा $b \in \Domain{M_2}$ घ्या.
  \end{enumerate}
  $\Struct{M_2}$ हा अंतबिंदुविरहित सघन रेषीय क्रम असल्यामुळे अपेक्षित
  गुणधर्म असलेला $b$ नेहमी शोधता येतो. $q = p \cup \{ \langle a, b
  \rangle \}$ असे परिभाषित करा; मग $q \in \PIso{I}$ हा $p$ चा अपेक्षित
  विस्तार आहे. यावरून पुढे गुणधर्म प्रस्थापित होतो. मागे गुणधर्मही
  याचप्रमाणे सिद्ध होतो. म्हणून $\Struct{M_1} \iso[p]
  \Struct{M_2}$; \olref[pis]{thm:p-isom1} नुसार $\Struct{M_1} \iso
  \Struct{M_2}$.
\end{proof}

\begin{prob}
  $\PIso{I}$ हा मागे गुणधर्म पूर्ण करतो, हे तपासून
  \olref[mod][bas][dlo]{thm:cantorQ} चा पुरावा पूर्ण करा.
\end{prob}

\begin{rem}
  $\Struct{S}$ हा कोणताही !!{enumerable} अंतबिंदुविरहित सघन रेषीय क्रम
  असू द्या. मग (\olref{thm:cantorQ} नुसार) $\Struct{S} \iso
  \Struct{Q}$, जेथे $\Struct{Q} = (\Rat, <)$ हा परिमेय संख्यांचा संच
  $\Rat$ प्रांत म्हणून असलेला !!{enumerable} सघन रेषीय क्रम आहे. आता
  \olref[thm]{remark:R} मधील !!{structure}~$\Struct{R} = (\Real, <)$
  पुन्हा विचारात घ्या. असा !!a{enumerable} !!{structure}~$\Struct{S}$
  आहे की $\Struct{R} \elemequiv \Struct{S}$, हे आपण पाहिले. पण
  $\Struct{S}$ हा !!a{enumerable} अंतबिंदुविरहित सघन रेषीय क्रम आहे;
  म्हणून तो !!{structure}~$\Struct{Q}$ शी समरूपी (आणि म्हणून प्राथमिक
  सममूल्य) आहे. प्राथमिक सममूल्यतेच्या संक्रमणकतेने $\Struct{R}
  \elemequiv \Struct{Q}$. (हाच पुढे-मागे युक्तिवाद वापरून
  $\Struct{R} \iso[p] \Struct{Q}$ प्रस्थापित करून आपण हे थेटही दाखवू
  शकलो असतो.)
\end{rem}
\end{document}
